\clearpage
\section{Literaturrecherche}

\subsection{Notizen zu Papers}
\textbf{Comparative analysis of real-time precise point positioning zenith total delay estimates}
ZTD ist von Wasserdampf in Atmosphäre abhängig. Kann rückwärts verwendet werden, um den Wasserdampf 
in der Atmosphäre zu bestimmen. Wird in Wettervorhersage verwendet. Um in schnell-aktualisierenden 
Wettermodellen ZTD verwenden zu können, müssen diese alle 5-10 mit Genauigkeit von 5-30 mm bestimmt werden. 
IGS Stellt real time orbit und clock Daten zur Verfügung (boradcast Ephemeriden, Orbit und Uhrkorrekturen). 
Dies ermöglicht PPP in Realtime zu berechnen, effizienter als Doppeldifferenzen.

In PPP ist ZTD anfällig für die radiale Komponente im Orbitfehler. Ionosphäre-Verzögerungen erster Ordnung 
werden durch 2-Frequenzen eliminiert. Es bleiben aber Restverzögerungen höhrerer Ordnung, welche im ZTD 
aufgehen. Diese sind linear abhängig von der Sonnenaktivität. Bei starker Sonnenaktivität steigt somit der 
ZTD-Fehler aus der Ionosphäre von 0.6 auf 4 mm. Weiter ist er vom Druck an der Erdoberfläche abhängig 
(-0.1 - -0.2 mm/hPa). Anntenenfehler (Phase center offset, phase center variations, daome geometry) haben auch 
Einfluss auf ZTD v0n 2-10 mm. Tropospheric mapping functions haben auch Einfluss auf ZTD. Dieser ist kleiner, 
je grösser der cut-off Winkel.

Je nach Software-Lösung zur Berechnung der PPP ZTD unterschiedliche Genauigkeiten. Es ist erreichbar mit 
gewissen Lösungen. 

\textbf{Pushing the sensitivity limits of RTS-based continuous deformation monitoring of an alpine valley}
Nichts zu ZTD drin

\textbf{tropospheric delay parameters derived from gnss-tracking data of a fast moving fleet of train}
Aufbau Atmosphäre Trophosphärischer Delay: Aufteilbar in hydrostatischer Anteil (Luftdruck) und feuchter 
Anteil (Wasserdampf).  Weitere Unterscheidung nach Richtung (Zenith Delay / Slant delay (in line of sight 
direvtion to satellite)

ZTD Verzögerung in Zenithrichtung aufgrund Troposphäre. Trockener Anteil deutlich grösser als feuchter 
Anteil. Feuchert Anteil dafür sehr lokal und zeitlich variabel, trockener Anteil besser modellierbar 
(z.B. Saastamoinen)

Slant delay: Zenithdelay mit mapping functions auf die Elevation hochgerechnet. Vienna Mapping Function 
aktuell als genaueste Mapping Function angesehen. Auf Relation von hydrostatischer Mappingfunktion zusammen 
mit Höhe. Aktuell VMF3 neuste Version, insbesondere in tiefen Winkeln besser

Viele Wettermodelle können heute ZTD oder Produkte die daraus abgleitet werden einberechnen. Dabei können 
entweder Einzelmessungen verwendet werden oder GNSS Tomography gemacht werden. Hier wird über ein Referenznetz 
alle möglichen ZTD berechnet. Daraus wird ein 3D Feld aus Voxeln berechnet, welche die Atmosphäre darstellt 
berechnet. 

GNSS Radio Occultation ist Space-based Technik zur Bestimmung von Atmosphärischen Parametern. Hierzu werden 
GNSS-Signale von LEO Satelliten empfangen. Daraus lässt sich Laufzeitverzögerung und Biegung durch die 
Troposphäre berechnen. Gute vertikale Auflösung, schlechte Horizontale Auflösung.
Schätzen von Troposphärischen Parametern aus GNSS-Auswertung: Parameter können in GNSS-Berechnung geschätzt 
werden. Hierbei ist Höhenfehler und ZTD stark korreliert. Je niedriger der Cut-off angle umso besser können 
sie auseinandergehalten werden. Wird ZTD in Auswertung geschätzt
